\documentclass[11pt]{article}
\usepackage[margin=25mm]{geometry}
\usepackage{fontspec}
\setmainfont{Latin Modern Roman}
\usepackage{amsmath,amssymb,mathtools}
\usepackage{unicode-math}
\setmathfont{Latin Modern Math}
\usepackage{xurl,hyperref}
\hypersetup{colorlinks=true,urlcolor=blue}
\setcounter{secnumdepth}{0}
\setlength{\parindent}{0pt}
\setlength{\parskip}{5pt}
\setlength{\emergencystretch}{3em}
\newcommand{\sref}[2]{\hyperlink{source:#1:#2}{[S#2]}}
\newcommand{\cataloguescope}[1]{\paragraph{Recorded target.}#1}
\newcommand{\R}{\mathbb R}
\newcommand{\C}{\mathbb C}
\newcommand{\N}{\mathbb N}
\newcommand{\Z}{\mathbb Z}
\newcommand{\Q}{\mathbb Q}
\newcommand{\F}{\mathbb F}
\newcommand{\eps}{\varepsilon}
\begin{document}
% UnsolvedMath ID 8000009; source catalogue code AMR-079-0009
\setcounter{section}{8000008}
\section{The parking problem}\label{8000009}
{\small UnsolvedMath ID 8000009 · Mathematical Physics}
\subsection{Definitions and mathematical statement}

\paragraph{Shared notation.}$\mathbb N=\{1,2,\ldots\}$; the standard number fields have their usual meanings. All additional conventions are stated locally.


\paragraph{Definitions and notation.}The target concerns observed spacings between real cars parked on a long street. A microscopic stochastic model specifies the actions and interactions of individual drivers and cars, from which spacing statistics are derived. GUE and GOE are Gaussian unitary and orthogonal random-matrix ensembles; their normalized local-spacing laws are different. Random sequential interval deposition is one possible parking model, not the model prescribed by this source.

\paragraph{Catalogue wording retained for traceability.}
Develop a microscopic model for parked-car spacings explaining the source’s observed random-matrix behavior: GUE statistics for two-way-street parking on the right and GOE statistics for left-side parking on one-way streets.
\paragraph{Scope and source qualification.}The source refers to empirical spacing studies by Šeba and others. A theorem for the classical random sequential adsorption model alone would be a different target unless it explains those observations.
\subsection{Short English statement}
Build a microscopic model explaining the different random-matrix spacing statistics observed for parked cars.
\subsection{Research attempt}
\paragraph{Model, date and investigated scope.}
OpenAI GPT-6 Astra (configured model \texttt{gpt-6-astra}), ultra reasoning; generated 11 October 2026. Derived the exact marginal and gamma limit of a constrained product-gap model, showing that matching repulsion exponents alone does not explain the observed GUE/GOE distinction.

\paragraph{Sources and literature check.}
Read the supplied complete definition, catalogue wording, scope qualifications and source list for this canonical ID. The original-source citations are inherited as supplied references; no fresh primary-source or novelty search was performed for this attempt. No current open-status claim is inferred from this calculation.

\paragraph{A local-gap model and what it can explain.}Consider \(N\ge2\) nonnegative gaps \(s_i\) with total free length \(L\), and postulate a density on the simplex proportional to \(\prod_i s_i^\beta\), where \(\beta>-1\). This is a modelling choice, not a derived driver law. Integrating out \(s_2,\ldots,s_N\) by scaling them with \(L-s_1\) gives a one-gap density
\[p_N(s)=C_N s^\beta(L-s)^{(N-1)(\beta+1)-1},\quad0<s<L,
\]
where \(C_N\) normalizes the integral. At fixed \(N\), the small-gap exponent is exactly \(\beta\), so choosing \(\beta=1\) or \(2\) can mimic different repulsion powers.
\paragraph{The large-system marginal.}Take \(L=N\), giving mean gap one. The law is equivalently obtained by independent gamma variables \(Y_i\) of shape \(\beta+1\) and common rate one, setting \(s_i=NY_i/\sum_jY_j\); the change of variables from total \(\sum Y_j\) and normalized proportions factors the gamma density. Since the sample mean tends to \(\beta+1\), each fixed \(s_i\) converges in distribution to \(Y_i/(\beta+1)\). Its density is
\[p(s)=\frac{(\beta+1)^{\beta+1}}{\Gamma(\beta+1)}s^\beta e^{-(\beta+1)s},\quad s>0.
\]
The mean and variance of the gamma variables are finite, so even Chebyshev's inequality suffices for the required convergence of their mean. This limiting law has an exponential tail; choosing a repulsion exponent has not selected a random-matrix spacing law.
\paragraph{First gap.}The gap-product ansatz inserts \(\beta\) by hand and gives no mechanism by which traffic direction or parking side changes it. More fundamentally, a local one-gap interaction need not reproduce the long-range correlations of a random-matrix point process. A microscopic driver-and-car rule leading to the observed two symmetry classes, with a check against empirical observations, remains missing. Random sequential deposition would likewise require a demonstrated connection to those observations. No current empirical claim or universal parking law is made here.

\paragraph{Verification and adversarial audit.}
Checked simplex exponent, total-length normalization, gamma representation, mean-one scaling and the unproved status of the physical model. This was an LLM adversarial audit, not independent human verification; no theorem prover was used. Prompts A then B applied to this ID: Checked simplex exponent, total-length normalization, gamma representation, mean-one scaling and the unproved status of the physical model. Scope remains limited: Derive a microscopic parking dynamics whose stationary correlations and traffic-direction dependence yield the specific observed random-matrix laws. Final pass clarified the nondegenerate simplex condition N>=2.

Coordinator checked the simplex marginal, gamma-variable representation and the scaled limiting density, including the distinction between repulsion exponent and a full spacing law. This additional check was model-assisted, not independent human verification.

\paragraph{Tools, process and contributions.}
Python 3 and local shell used to read supplied files and serialize this manuscript; mathematics derived symbolically in the text, with no numerical experiment or formal proof assistant. Contribution prompt A was followed by prompt B. GPT-6 Astra with ultra reasoning drafted the research and performed its stated self-audit; any additional coordinator contribution is identified in the verification and audit. A GPT-6-based Codex coordinator assembled the manuscript and checked repository integration; its exact variant and reasoning configuration were not exposed. No human mathematical contribution, independent human proof check or formal proof certification is supplied. Complete derivations and the consulted source references are included here; private reasoning transcripts are not published.

\paragraph{Final status and first remaining gap.}
\textbf{UNRESOLVED}. Derive a microscopic parking dynamics whose stationary correlations and traffic-direction dependence yield the specific observed random-matrix laws.

\section{Completion Estimate}
\noindent\textbf{COMPLETION: 0\%}
The calculation is an elementary benchmark or conditional reduction; it establishes no claimed new portion of the complete source target.

\subsection{Sources}
\begin{itemize}
\item[\hypertarget{source:8000009:1}{S1}] ULAM.AI and the UnsolvedMath Contributors, UnsolvedMath, record 8000009 (AMR-079-0009). Catalogue curation under CC BY 4.0; definitions and English exposition added locally.
\url{https://huggingface.co/datasets/ulamai/UnsolvedMath}.
\item[\hypertarget{source:8000009:2}{S2}] Percy Deift, Some open problems in random matrix theory and the theory of integrable systems (2007), arXiv:0712.0849, Problem 9.
\url{https://arxiv.org/abs/0712.0849}.
\end{itemize}
\label{end:8000009}
\end{document}
